
Throughout this appendix, $\calT$ denotes a \GGM\ seed tree with leaf set
$\{0,\ldots,t-1\}$ and $\calF : V(\calT) \to \{0,1\}$ its flag labelling,
as produced by the $\Publish$ phase of the relevant \GZKST\ instance.
We use the convention $\calF(v) = 1$ to mean that $v$'s subtree contains at
least one hidden leaf ($i \in \calH$), and $\calF(v) = 0$ to mean that all
leaf descendants of $v$ are revealed ($i \in \calR$).

% ------------------------------------------------------------
\subsection{Proof of Proposition~\ref{prop:less-gzkst}
  (LESS-v2 is a \GZKST\ instance)}
\label{app:proof-less}
% ------------------------------------------------------------

Let $\calT$ be the LESS-v2 seed tree with $\lceil \log_2 t \rceil$ levels.
After the three-phase flag-tree construction the flag values are:
\begin{align}
  \calF(\mathrm{leaf}_i) &= \mathbf{1}[i \in \calH], \quad i \in \{0,\ldots,t-1\},
  \label{eq:leaf-flag}\\
  \calF(v) &= \calF(\ell) \vee \calF(r), \quad \text{for every internal node } v
  \text{ with children } \ell, r.
  \label{eq:or-prop}
\end{align}
Equations~\eqref{eq:leaf-flag}--\eqref{eq:or-prop} jointly imply:
\begin{equation}
  \calF(v) = 1 \iff \exists\, i \in \calH : \mathrm{leaf}_i
  \text{ is a descendant of } v.
  \label{eq:flag-char}
\end{equation}
The published path is
$\tpath = \{v : \calF(v) = 0,\; \calF(\mathrm{parent}(v)) = 1\}$
(with the convention that $\calF(\mathrm{parent}(\mathrm{root})) = 1$).

\begin{proof}[Proof of Invariant~\ref{inv:correctness} (Correctness)]
Let $i \in \calR$.  Consider the root-to-leaf path
$\pi_i = (v_0 = \mathrm{root},\, v_1,\, \ldots,\, v_d = \mathrm{leaf}_i)$.

Since $|\calH| \geq 1$ (the challenge weight satisfies $w \geq 1$), at least one
leaf has flag~$1$.  By~\eqref{eq:flag-char}, $\calF(v_0) = \calF(\mathrm{root}) = 1$.

Since $i \in \calR$, equation~\eqref{eq:leaf-flag} gives $\calF(v_d) = 0$.

Because $\calF(v_0) = 1$ and $\calF(v_d) = 0$, the path $\pi_i$ contains
at least two nodes with different flags.  Define
$k = \min\{j \in \{0,\ldots,d\} : \calF(v_j) = 0\}$.
Then $k \geq 1$ (since $\calF(v_0) = 1$), $\calF(v_k) = 0$, and
$\calF(v_{k-1}) = 1$.  Hence $v_k \in \tpath$, and $v_k$ lies on
the root-to-leaf path $\pi_i$, establishing Invariant~\ref{inv:correctness}.
\end{proof}

\begin{proof}[Proof of Invariant~\ref{inv:hiding} (ZK hiding)]
Let $i \in \calH$ and let $v$ be any proper ancestor of $\mathrm{leaf}_i$
(including the root).  Since $\mathrm{leaf}_i$ is a descendant of $v$ and
$i \in \calH$, equation~\eqref{eq:flag-char} gives $\calF(v) = 1$.
Hence no ancestor of $\mathrm{leaf}_i$ can satisfy $\calF(v) = 0$, so no
ancestor of $\mathrm{leaf}_i$ lies in $\tpath$.  Invariant~\ref{inv:hiding}
holds.
\end{proof}

\begin{proof}[Extension to incomplete trees (Remark~\ref{rem:incomplete})]
When $t$ is not a power of~$2$, the LESS-v2 tree is incomplete: some leaves
appear at depth $\lceil \log_2 t \rceil - 1$ rather than $\lceil \log_2 t \rceil$.
The index-tracking arrays $\mathsf{npl}$, $\mathsf{ll}$, $\mathsf{off}$, and
$\mathsf{lsi}$ ensure that $\LabelLeaves$ correctly identifies which nodes are
leaves and assigns flags according to~\eqref{eq:leaf-flag}, and that
$\ComputeSeeds$ performs the bottom-up OR pass~\eqref{eq:or-prop} over the
actual tree structure.

The proofs of both invariants above are purely graph-theoretic and depend only
on the properties~\eqref{eq:leaf-flag}--\eqref{eq:flag-char}, which hold
regardless of whether leaves appear at a uniform depth.  Hence both invariants
hold for any incomplete tree shape produced by the LESS-v2 expansion.
\end{proof}

% ------------------------------------------------------------
\subsection{ Proof of Proposition~\ref{prop:meds-gzkst}
  (MEDS is a \GZKST\ instance)}
\label{app:proof-meds}
% ------------------------------------------------------------

MEDS uses $\SeedTreeToPath$, which implements a top-down erasure: starting from
a fully built tree, for each $i \in \calH$ the algorithm removes the label of
$\mathrm{leaf}_i$ and propagates the removal upward, erasing every ancestor
whose entire subtree has had all labels removed.

\begin{claim}
After the erasure pass, a node $v$ retains its label if and only if
$\calF(v) = 0$ in the sense of~\eqref{eq:flag-char}, i.e., $v$'s subtree
contains no hidden leaf.
\end{claim}
\begin{proof}
A node is erased iff every leaf in its subtree has been removed, which holds
iff every leaf in its subtree is in $\calH$.  Equivalently, a node retains its
label iff its subtree contains at least one revealed leaf, i.e., $\calF(v) = 0$.
\end{proof}

The published path $\tpath$ consists of the highest surviving nodes in the tree,
serialised in depth-first order; equivalently,
$\tpath = \{v : v \text{ retains its label},\; \mathrm{parent}(v) \text{ was erased}\}$.
By the claim, this is $\{v : \calF(v) = 0,\; \calF(\mathrm{parent}(v)) = 1\}$,
which is identical to the LESS-v2 definition.

Invariants~\ref{inv:correctness} and~\ref{inv:hiding} follow by the same
argument as in Appendix~A.1.

% ------------------------------------------------------------
\subsection{Proof of Proposition~\ref{prop:cross-gzkst}
  (CROSS is a \GZKST\ instance)}
\label{app:proof-cross}
% ------------------------------------------------------------

The seed-tree component of CROSS ($\SeedTreePaths$) is structurally equivalent
to MEDS's $\SeedTreeToPath$ (see Section~\ref{sec:mapping}).  The argument of
Appendix~A.2 therefore applies directly, establishing
Invariants~\ref{inv:correctness} and~\ref{inv:hiding} for the seed-hiding part.

The additional Merkle proof over $\{\mathsf{cmt}_0[i]\}_{i \in \calH}$ is an
authenticity mechanism for the commitment digests that the verifier cannot
recompute from the revealed seeds.  It neither publishes any seed in
$\{\eseed_i : i \in \calH\}$ nor affects the seed-path computation, and is
therefore orthogonal to both invariants.

% ------------------------------------------------------------
\subsection{Proof of Proposition~\ref{prop:mpcith-gzkst}
  (MPCitH/VOLEitH family is a \GZKST\ instance)}
\label{app:proof-mpcith}
% ------------------------------------------------------------

For each epoch $e \in \{1,\ldots,\tau\}$, the tree has $N_e$ leaves with
exactly one hidden leaf at position $i^*[e]$.
$\mathsf{GetSiblingPath}$ publishes the \emph{co-path}: for each node on the
root-to-hidden-leaf path, the sibling is added to $\tpath$.

Formally, let $\pi^*_e = (u_0 = \mathrm{root},\, u_1,\, \ldots,\, u_d = \mathrm{leaf}_{i^*[e]})$
be the root-to-hidden-leaf path in epoch $e$.  The co-path for epoch $e$ is
$\tpath_e = \{ \mathrm{sibling}(u_k) : k = 1, \ldots, d \}$,
and $\tpath = \bigcup_e \tpath_e$.

\begin{proof}[Proof of Invariant~\ref{inv:correctness}]
Let $(e, i)$ with $i \neq i^*[e]$ be a revealed pair.  Let $k$ be the shallowest
level at which the path from the root to $\mathrm{leaf}_i$ and the path
$\pi^*_e$ diverge; denote the node on $\pi^*_e$ at level $k$ by $u_k$ and its
sibling by $\mathrm{sibling}(u_k)$.

Since the paths diverge at level $k$, the subtree below $\mathrm{sibling}(u_k)$
contains $\mathrm{leaf}_i$ but not $\mathrm{leaf}_{i^*[e]}$.
Since $\mathrm{sibling}(u_k) \in \tpath_e$, the verifier holds this node and
can expand its subtree downward to recover $\eseed_{(e,i)}$.
\end{proof}

\begin{proof}[Proof of Invariant~\ref{inv:hiding}]
The co-path $\tpath_e$ consists of siblings of nodes on $\pi^*_e$, never
nodes on $\pi^*_e$ itself.  In particular, no ancestor of $\mathrm{leaf}_{i^*[e]}$
is in $\tpath_e$.  Expanding any node in $\tpath_e$ downward therefore never
reaches $\mathrm{leaf}_{i^*[e]}$.

For epoch $e' \neq e$, the co-paths $\tpath_{e'}$ belong to a different
epoch's tree (or a different sub-tree of the interleaved structure), and
contain no node that is an ancestor of $\mathrm{leaf}_{i^*[e]}$.
Hence $\tpath$ contains no ancestor of any hidden leaf, establishing
Invariant~\ref{inv:hiding}.
\end{proof}

% ------------------------------------------------------------
\subsection{Proof of Proposition~\ref{prop:alg-correctness}
  (Correctness of Algorithm~\ref{alg:key-recovery})}
\label{app:proof-recovery}
% ------------------------------------------------------------

Let $m = |\calS| = s - 1$.  We model collection of secret components as a
\emph{generalised coupon-collector} problem with $m$ distinct coupons.

\paragraph{Correctness of Extract.}
For each scheme, Definition~\ref{def:extract} specifies $\Extract$ as a closed-form
algebraic procedure: matrix inversion over $\mathbb{F}_q$ (MEDS), column extraction
(LESS), a group operation (CROSS), or linear combination of shares (MPCitH).
Each procedure is deterministic and correct given the dual revelation
$(\eseed_i, \resp_i)$: the algebraic relationship between the response and the
hidden seed is an identity (not an approximation), so $\Extract$ returns exactly
$\sk' = $ the secret component indexed by $j = c_i$.

\paragraph{Termination.}
Define the \emph{progress probability} at step $q$ as
\[
  p_{\min} \;=\; 1 - \left(1 - \frac{1}{s-1}\right)^{\!\bar{w}},
\]
where $\bar{w} = w \cdot \ell / t \geq 1$ is the expected number of hidden
leaves in the subtree of $v^*$.  Since each hidden leaf's challenge index
$c_i$ is uniform in $\{1,\ldots,s-1\}$ independently of which components
have already been collected, we have
\[
  \Pr\bigl[\mathsf{new}(q) \geq 1 \;\big|\; \mathsf{Collected}\bigr]
  \;\geq\; p_{\min} \;>\; 0
  \qquad \text{for all } s \geq 2,\; \bar{w} \geq 1.
\]
Let $T_j$ be the number of effective queries to go from
$|\mathsf{Collected}| = j-1$ to $|\mathsf{Collected}| = j$.
Conditional on $|\mathsf{Collected}| = j-1$, each query succeeds with
probability at least $p_{\min}$, so $T_j$ is stochastically dominated by
a $\mathrm{Geom}(p_{\min})$ random variable, which is finite almost
surely.  Hence $Q = \sum_{j=1}^{m} T_j$ is finite almost surely, and
Algorithm~\ref{alg:key-recovery} terminates with probability~$1$.

\paragraph{Expected query count and tail bound.}
Let $R_k$ be the number of uncollected secret indices after $k$ effective
queries, and let $r = R_{k-1} > 0$.  Each of the $\bar{w}$ hidden leaves
in the subtree independently draws a uniform index from $\{1,\ldots,s-1\}$,
so the probability that \emph{none} of the $\bar{w}$ draws falls in the
remaining set of size $r$ is $\bigl((s-1-r)/(s-1)\bigr)^{\bar{w}}$.
Hence the probability of revealing at least one new index is
\begin{equation}
  p_r \;=\; 1 - \left(\frac{s-1-r}{s-1}\right)^{\!\bar{w}}
           \;=\; 1 - \left(1 - \frac{r}{s-1}\right)^{\!\bar{w}}.
  \label{eq:pr-exact}
\end{equation}
Since $r \geq 1$ and $\bigl(1 - r/(s-1)\bigr)^{\bar{w}} \leq
\bigl(1 - 1/(s-1)\bigr)^{\bar{w}}$, we have $p_r \geq p_{\min}$ for
all $r \geq 1$.

The expected total number of effective queries is
\begin{equation}
  \mathbb{E}[Q] \;=\; \sum_{r=1}^{m} \frac{1}{p_r}
  \;\leq\; \frac{m}{p_{\min}} \;=\; \frac{s-1}{p_{\min}},
  \label{eq:exp-Q}
\end{equation}
where the equality uses the standard coupon-collector identity for the
case where each query attempts $\bar{w}$ independent draws.

For the tail bound, each $T_j \leq k_0$ with probability at least
$1 - (1 - p_{\min})^{k_0}$ (geometric tail), so by a union bound over
$m = s-1$ stages,
\[
  \Pr[Q > k] \;\leq\; (s-1)\cdot(1-p_{\min})^{\lfloor k/(s-1)\rfloor}.
\]
Setting $(s-1)(1-p_{\min})^{\lfloor k/(s-1)\rfloor} \leq \delta$ and
solving gives $k \leq \bigl\lceil \tfrac{s-1}{p_{\min}} \cdot
\ln\tfrac{s-1}{\delta} \bigr\rceil$, establishing the bound in
Proposition~\ref{prop:alg-correctness}.

\paragraph{Output correctness.}
When the loop terminates, $|\mathsf{Collected}| = m$, and for each
$j \in \{1,\ldots,s-1\}$ there is exactly one entry $(j, \sk'_j) \in \mathsf{Collected}$
with $\sk'_j = A_j^{-1}$ (or the corresponding secret for each scheme).
The $\mathsf{Collect}$ step assembles these into $\widehat{\sk}$, which equals $\sk$
by correctness of $\Extract$.

% ------------------------------------------------------------
\subsection{Derivation of Proposition~\ref{prop:query-count}
  (Expected query count)}
\label{app:proof-query}
% ------------------------------------------------------------

\paragraph{Setup.}
Let $v^*$ be the faulted node with $\ell$ leaf descendants.
By the Fiat--Shamir uniformity assumption, the challenge index $c_i$ is
uniform in $\{1,\ldots,s-1\}$ independently for each
$i \in \calH \cap \mathrm{leaves}(v^*)$.
The actual number of hidden leaves in the subtree,
$L = |\calH \cap \mathrm{leaves}(v^*)|$, is a hypergeometric random variable
with mean $\bar{w} = w \cdot \ell / t$.

\paragraph{Distinct secrets per query (fixed $L = \ell_h$).}
Conditioning on $L = \ell_h$, the number of \emph{distinct} secret indices
revealed in one effective query has expectation
\begin{align}
  \mathbb{E}[\text{\#distinct} \mid L = \ell_h]
  &= (s-1)\!\left(1 - \left(1 - \frac{1}{s-1}\right)^{\!\ell_h}\right),
  \label{eq:edist-exact}
\end{align}
by linearity of expectation and symmetry of the uniform distribution.

\paragraph{Averaging over $L$ and the Jensen remark.}
Taking expectation over $L$,
\[
  \mathbb{E}[\text{\#distinct}]
  = (s-1)\!\left(1 - \mathbb{E}\!\left[\left(1-\tfrac{1}{s-1}\right)^{\!L}\right]\right).
\]
Since $f(x) = \alpha^x$ with $\alpha = 1 - 1/(s-1) \in (0,1)$ satisfies
$f''(x) = (\ln\alpha)^2 \alpha^x > 0$ (convex), Jensen's inequality gives
$\mathbb{E}[f(L)] \geq f(\mathbb{E}[L]) = \alpha^{\bar{w}}$.
Hence
\begin{equation}
  \mathbb{E}[\text{\#distinct}]
  \;\leq\; (s-1)\!\left(1 - \left(1-\frac{1}{s-1}\right)^{\!\bar{w}}\right)
  \;=:\; E_{\mathrm{dist}}.
  \label{eq:edist}
\end{equation}
That is, $E_{\mathrm{dist}}$ is an \emph{upper bound} on the true expected
number of distinct secrets per query; the exact value requires knowing the
distribution of $L$.  For a deterministic fault on a fixed node this
distinction vanishes ($L$ is determined by the challenge), so
$E_{\mathrm{dist}}$ is exact when conditioned on a specific query.

\paragraph{Total queries.}
From equation~\eqref{eq:pr-exact} in Appendix~A.5,
$p_r = 1 - \bigl((s-1-r)/(s-1)\bigr)^{\bar{w}}$,
and the coupon-collector identity gives
\[
  \mathbb{E}[Q] = \sum_{r=1}^{m} \frac{1}{p_r} \leq \frac{m}{p_{\min}}
  = \frac{s-1}{p_{\min}},
\]
since $p_r \geq p_{\min}$ for all $r \geq 1$.  Note that
$p_{\min} = E_{\mathrm{dist}}/(s-1)$, so this bound equals
$(s-1)^2/E_{\mathrm{dist}}$, which for $p_{\min} \approx 1$ (large~$w$)
gives $\mathbb{E}[Q] \lesssim s-1$.